\documentclass[12pt,a4paper]{article}
\usepackage[T2A]{fontenc}
\usepackage[utf8]{inputenc}
\usepackage[russian]{babel}
\usepackage{cmap}
\usepackage{amsmath,amssymb,mathtools,geometry,enumitem}
\geometry{margin=2cm}
\setlength{\parindent}{0pt}
\setlist[enumerate]{leftmargin=*}

\begin{document}

\section*{Рабочая тетрадь: квадратные уравнения}

\subsection*{Теория. Дискриминант}
Для квадратного уравнения
\[
ax^2 + bx + c = 0, \qquad a \neq 0
\]
дискриминант вычисляется по формуле
\[
D = b^2 - 4ac.
\]

Если:
\begin{itemize}
\item $D > 0$, уравнение имеет два различных корня;
\item $D = 0$, уравнение имеет один корень;
\item $D < 0$, действительных корней нет.
\end{itemize}

\subsection*{Задачи}
\begin{enumerate}
\item \textbf{Найдите корни уравнения}
\[
x^2 - 5x + 6 = 0.
\]

\vspace{6mm}

\item \textbf{Найдите корни уравнения}
\[
x^2 - 9 = 0.
\]

\vspace{6mm}

\item \textbf{Определите число действительных корней}
\[
x^2 + 4x + 5 = 0.
\]
\end{enumerate}

\end{document}
